\begin{resumo}[Abstract]
 \begin{otherlanguage*}{english}
Brazilian federal agencies monitor public policies with data spread across
systems that were built and are managed separately, and few of them share a
key a join can use directly. Their identifiers differ in format, must be
combined, or appear only inside free text. In Gov Hub, an open-source data
platform built by the University of Brasília with federal ministries, only
one of 29 cross-source join conditions in production code is a direct
equality between untransformed columns. Schema matchers do not produce these
transformations, and a Large Language Model (LLM) that writes the join SQL
can return a wrong join that still runs. This work proposes a system in which
the LLM only chooses the key, from evidence collected beforehand and a closed
catalog of operations, or abstains, while deterministic code generates and
checks the join in dbt. The same principle, which we call bounded agency,
applies to the coding agent that builds the system through Spec-Driven
Development. The method is Design Science Research with a case study in the
Ministry of Culture's Gov Hub repository. The planned evaluation compares the
system with classical matchers, rules, and other LLM uses on a benchmark of
real and perturbed pairs, compares open and proprietary models, measures the
effect on analysts, applies the system to source pairs the ministry does not
yet join, and compares the agent's work with and without a spec using tests
hidden from it. This first stage presents the problem, the system design, the
study protocol, and preliminary results.

   \vspace{\onelineskip}

   \noindent
   \textbf{Keywords}: data integration. language models. spec-driven
   development. open government data. schema matching.
 \end{otherlanguage*}
\end{resumo}
